%% ============================================================
%%  TFM Paper — Laura González Villar
%%  Structure: Introduction · Related Work · Material · Methods
%%             · Experimental Results · Discussion · Conclusions
%%  Paste this file into Overleaf as main.tex
%% ============================================================

\documentclass[a4paper, 11pt]{article}

%% ── Encoding & language ──────────────────────────────────────
\usepackage[utf8]{inputenc}
\usepackage[T1]{fontenc}
\usepackage[english]{babel}

%% ── Layout ───────────────────────────────────────────────────
\usepackage[top=2.5cm, bottom=2.5cm, left=3cm, right=3cm]{geometry}
\usepackage{setspace}
\onehalfspacing

%% ── Fonts & typography ───────────────────────────────────────
\usepackage{lmodern}
\usepackage{microtype}

%% ── Math ─────────────────────────────────────────────────────
\usepackage{amsmath, amssymb}

%% ── Tables ───────────────────────────────────────────────────
\usepackage{booktabs}
\usepackage{multirow}
\usepackage{array}

%% ── Figures ──────────────────────────────────────────────────
\usepackage{graphicx}
\usepackage{float}

%% ── Colors & listings ────────────────────────────────────────
\usepackage{xcolor}
\usepackage{listings}

%% ── Hyperlinks ───────────────────────────────────────────────
\usepackage[colorlinks=true, linkcolor=blue, citecolor=blue, urlcolor=blue]{hyperref}

%% ── Bibliography ─────────────────────────────────────────────
\usepackage[numbers, sort&compress]{natbib}

%% ── Algorithm ────────────────────────────────────────────────
\usepackage{algorithm}
\usepackage{algpseudocode}

%% ── Custom commands ──────────────────────────────────────────
\newcommand{\eg}{\textit{e.g.,}\ }
\newcommand{\ie}{\textit{i.e.,}\ }
\newcommand{\etal}{\textit{et al.}\ }


%% ============================================================
%%  TITLE
%% ============================================================
\title{%
  \textbf{Explainable ADME Property Prediction Using LightGBM,\\
  SHAP, and Large Language Models}
  \vspace{0.3em}
}

\author{%
  Laura González Villar\\
  \small Màster en Ciència de Dades (MDS), Facultat d'Informàtica de Barcelona\\
  \small Universitat Politècnica de Catalunya (UPC)\\
  \small \href{mailto:lglezvi@gmail.com}{lglezvi@gmail.com}
  \and
  % Add supervisors / co-authors here if required by journal
  % Albert Vellido \quad Caroline König\\
  % \small Departament de Ciències de la Computació, UPC
}

\date{August 2026}


%% ============================================================
\begin{document}
\maketitle
\thispagestyle{empty}

%% ============================================================
%%  ABSTRACT
%% ============================================================
\begin{abstract}
Predicting Absorption, Distribution, Metabolism, and Excretion (ADME)
properties early in drug discovery is critical to reducing attrition rates
in clinical development. Machine learning models can provide fast,
cost-effective predictions, but their black-box nature hinders adoption in
regulatory-sensitive environments. In this work we present an end-to-end
explainable pipeline that combines LightGBM regressors, SHAP
(SHapley Additive exPlanations) feature attribution, and a large language
model (LLM) to produce human-readable explanations of individual ADME
predictions. We train six independent models on a publicly available
dataset of 3{,}521 drug-like compounds covering six pharmacokinetic
endpoints: human and rat liver microsomal clearance (HLM, RLM), human and
rat plasma protein binding (hPPB, rPPB), MDR1 efflux ratio (MDR1-ER), and
aqueous solubility (Sol). Models achieve Pearson correlation coefficients
between 0.554 and 0.749. Three prompt-engineering strategies (\textit{basic},
\textit{biomedical}, and \textit{reasoning}) are compared across different
numbers of SHAP features; the biomedical strategy with $k{=}3$ features
yields the most domain-coherent explanations. An LLM-as-a-Judge evaluation
framework using Claude Haiku 4.5 (Anthropic) as an independent judge shows
that GPT-4.1-mini explanations are preferred in 96.6\% of pairwise comparisons. The pipeline is deployed as an interactive web
application on Hugging Face Spaces, making explainable ADME prediction
accessible to non-expert users.

\medskip
\noindent\textbf{Keywords:} ADME prediction, explainability, LightGBM,
SHAP, large language models, drug discovery, pharmacokinetics, XAI.
\end{abstract}

\newpage
\tableofcontents
\newpage


%% ============================================================
%%  1. INTRODUCTION
%% ============================================================
\section{Introduction}
\label{sec:intro}

Drug attrition due to poor pharmacokinetic properties accounts for a
substantial fraction of late-stage clinical failures, highlighting the need
for reliable early-stage ADME profiling~\cite{hughes2011physiochemical}.
Experimental assays, while accurate, are time-consuming and
expensive~\cite{doak2014direct}. Computational prediction offers a
complementary, high-throughput approach that can guide compound
prioritisation before resources are committed to wet-lab work.

Machine learning (ML) methods have demonstrated competitive predictive
performance on ADME endpoints~\cite{fang2023prospective,
konig2024adme}. However, regulatory frameworks such as the EU AI
Act~\cite{euaiact2024} and the EU MDR~\cite{eumdr2017} increasingly require
that AI-driven decisions in high-stakes domains (\eg life sciences,
medicine) be explainable and auditable. Black-box models alone therefore
fall short of practical deployment requirements.

Explainability techniques—particularly SHAP~\cite{lundberg2017unified}—can
surface which molecular features drive a given prediction, but the resulting
numerical feature importance scores remain inaccessible to medicinal
chemists unfamiliar with ML. Recent advances in large language models
(LLMs) open the possibility of translating these scores into natural
language that domain experts can act upon~\cite{kojima2022large,
wei2022chain}.

\medskip
\noindent\textbf{Contributions.} This paper makes the following
contributions:
\begin{enumerate}
  \item We train and validate six LightGBM regressors for ADME endpoint
        prediction using RDKit physicochemical descriptors, establishing
        competitive baselines on the Fang \etal benchmark dataset.
  \item We integrate SHAP TreeExplainer to provide both global
        (dataset-level) and local (molecule-level) feature attributions.
  \item We design and evaluate three prompt-engineering strategies that
        translate SHAP values into biomedically coherent natural language
        explanations via GPT-4.1-mini.
  \item We evaluate explanation quality through an LLM-as-a-Judge
        pairwise framework and demonstrate the superiority of the
        \textit{biomedical} prompt strategy.
  \item We deploy the full pipeline as an open-access interactive web
        application.
\end{enumerate}

The remainder of the paper is organised as follows.
Section~\ref{sec:related} reviews related work.
Section~\ref{sec:material} describes the dataset and molecular
representation. Section~\ref{sec:methods} details the predictive and
explainability pipeline. Section~\ref{sec:results} presents experimental
results. Section~\ref{sec:discussion} discusses findings and limitations.
Section~\ref{sec:conclusions} concludes the paper.


%% ============================================================
%%  2. RELATED WORK
%% ============================================================
\section{Related Work}
\label{sec:related}

\subsection{ADME Property Prediction with Machine Learning}

Early ADME prediction models relied on physicochemical descriptors and
simple regression or classification methods~\cite{lipinski2001experimental}.
Gradient-boosted trees have since become competitive baselines owing to
their robustness to noisy biological data and their native interpretability
through feature importance~\cite{chen2016xgboost}. Fang~\etal~\cite{fang2023prospective}
compiled a large, carefully curated benchmark spanning six ADME endpoints
and demonstrated that LightGBM with RDKit descriptors achieves state-of-the-art
performance among ML approaches. König and Vellido~\cite{konig2024adme}
extended this line of work by applying XAI methods to the same benchmark,
establishing reference SHAP rankings that we use to validate our
explanations in Section~\ref{sec:results}.

\subsection{Explainability in Cheminformatics}

SHAP~\cite{lundberg2017unified} offers a theoretically grounded approach to
feature attribution based on Shapley values from cooperative game theory.
TreeExplainer~\cite{lundberg2020local} exploits the recursive structure of
decision trees to compute exact Shapley values in polynomial time, making it
practical for the large ensembles produced by gradient-boosted models.
SHAP has been applied to a range of cheminformatics tasks, including QSAR
modelling~\cite{rodriguez2021shap}, toxicity prediction, and property
prediction~\cite{konig2024adme}.

\subsection{Large Language Models for Scientific Explanation}

LLMs have demonstrated remarkable in-context reasoning capabilities
across scientific domains~\cite{kojima2022large, wei2022chain}.
Several recent works explore LLM-generated explanations as a bridge between
ML model outputs and domain-expert reasoning~\cite{bhatt2020explainable}.
Prompt engineering strategies—ranging from zero-shot to chain-of-thought
prompting~\cite{wei2022chain}—significantly affect the quality and domain
coherence of generated text. To our knowledge, this is the first work to
systematically study prompt strategies for translating SHAP values into
ADME-specific natural language explanations.

\subsection{LLM-as-a-Judge Evaluation}

Automated evaluation of free-form text quality using LLMs as judges was
formalised by Zheng~\etal~\cite{zheng2023judging}. The authors identified
\textit{self-preference bias}—the tendency of a model to rate outputs from
the same model family more favourably—as a key confound. We acknowledge
this limitation in our evaluation design (Section~\ref{sec:discussion}).


%% ============================================================
%%  3. MATERIAL
%% ============================================================
\section{Material}
\label{sec:material}

\subsection{Dataset}

We use the public ADME benchmark introduced by Fang~\etal~\cite{fang2023prospective},
which aggregates experimental measurements for drug-like small molecules
across six pharmacokinetic endpoints (Table~\ref{tab:endpoints}).
After filtering for valid SMILES and removing duplicates, the working
dataset contains \textbf{3{,}521 compounds}. All compounds are shared
across all six endpoints (\ie each molecule contributes to every model's
training set).

\begin{table}[H]
  \centering
  \caption{ADME endpoints, units, and dataset statistics.}
  \label{tab:endpoints}
  \begin{tabular}{llll}
    \toprule
    \textbf{Endpoint} & \textbf{Property} & \textbf{Unit} & \textbf{$n$} \\
    \midrule
    HLM      & Human liver microsomal clearance   & mL/min/kg & 3{,}521 \\
    RLM      & Rat liver microsomal clearance     & mL/min/kg & 3{,}521 \\
    hPPB     & Human plasma protein binding       & \%        & 3{,}521 \\
    rPPB     & Rat plasma protein binding         & \%        & 3{,}521 \\
    MDR1-ER  & MDR1 efflux ratio (P-gp)           & ratio     & 3{,}521 \\
    Sol      & Aqueous solubility                 & $\mu$g/mL & 3{,}521 \\
    \bottomrule
  \end{tabular}
\end{table}

Continuous endpoints (HLM, RLM, MDR1-ER, Sol) were log-transformed prior
to modelling to reduce skewness. PPB endpoints were kept on their original
percentage scale.

\subsection{Molecular Representation}

Each compound was represented as a vector of \textbf{217 physicochemical
descriptors} computed with RDKit~\cite{landrum2006rdkit}. Descriptors
cover lipophilicity (MolLogP), molecular weight, hydrogen-bond donors and
acceptors, topological polar surface area (TPSA), rotatable bonds,
aromaticity indices, ring counts, and various graph-theoretic topological
descriptors. Feature selection was not applied; all 217 descriptors were
used as input to the LightGBM models. This descriptor count differs
slightly from the 316 reported by König and Vellido~\cite{konig2024adme}
due to a difference in RDKit version.


%% ============================================================
%%  4. METHODS
%% ============================================================
\section{Methods}
\label{sec:methods}

\subsection{Predictive Models}

We trained six independent \textbf{LightGBM}~\cite{ke2017lightgbm}
regressors, one per ADME endpoint. LightGBM is a gradient-boosted decision
tree framework that employs histogram-based learning and leaf-wise growth,
offering fast training and competitive performance on tabular data.
Hyperparameters were tuned via 5-fold cross-validation using Optuna
(learning rate, number of leaves, minimum child samples, feature fraction,
bagging fraction). Final models were trained on the full dataset; the
same 80/20 train-test split was used consistently across all endpoints to
allow fair comparison.

\subsection{SHAP Feature Attribution}

Model explanations were computed using
SHAP~\cite{lundberg2017unified} with the \textbf{TreeExplainer}
backend~\cite{lundberg2020local}, which computes exact Shapley values
for tree-based models in polynomial time by exploiting the recursive
structure of decision trees.

\paragraph{Global explanations.}
For each endpoint we compute $|\phi_i|$ (mean absolute SHAP value) for
every descriptor across the full dataset, yielding a global feature
importance ranking that can be compared to the literature.

\paragraph{Local explanations.}
For a query molecule $m$, the top-$k$ features by $|\phi_i(m)|$ are
extracted as input to the LLM explanation stage (Section~\ref{subsec:llm}).
The SHAP baseline corresponds to the mean of model predictions on the
training set for the respective endpoint.

\subsection{LLM Explanation Pipeline}
\label{subsec:llm}

\subsubsection{Prompt Engineering Strategies}

Given the top-$k$ SHAP features and their signed values for a query
molecule, we generate a natural language explanation using
GPT-4.1-mini via the OpenAI API. Three prompt strategies were
designed and compared:

\begin{enumerate}
  \item \textbf{Basic}: A minimal prompt providing feature names and SHAP
        values with no domain framing.
  \item \textbf{Biomedical}: A structured prompt that contextualises the
        SHAP values within pharmacokinetic and medicinal chemistry
        knowledge, explicitly instructing the model to avoid ML
        terminology.
  \item \textbf{Reasoning}: A chain-of-thought prompt~\cite{wei2022chain}
        that asks the model to reason step-by-step before generating the
        explanation.
\end{enumerate}

Each strategy was evaluated for $k \in \{2, 3, 5\}$ top SHAP features,
yielding $3 \times 3 = 9$ experimental conditions. All explanations follow
a standardised four-part structure: (1) endpoint definition and importance,
(2) model prediction with uncertainty, (3) driver features and their
pharmacokinetic relevance, and (4) actionable medicinal chemistry
suggestions.

\subsubsection{Models Compared}

The LLM evaluation compares six models across two families:
GPT-4.1-mini (primary), GPT-4.1-nano, GPT-4o-mini (OpenAI); and
LLaMA-3-8B, Mistral-7B, Gemma-7B (open-source, via Hugging Face).

\subsection{LLM-as-a-Judge Evaluation}

Explanation quality was assessed via \textbf{pairwise comparisons}
following Zheng~\etal~\cite{zheng2023judging}. For each evaluation sample,
an independent LLM judge is shown two anonymised explanations generated by
different models and asked to select the better one based on four criteria:
domain accuracy, biomedical coherence, absence of ML jargon, and
actionability.

To mitigate self-preference bias~\cite{zheng2023judging}---the tendency of
LLM judges to favour outputs from models of the same family---we use
\textbf{Claude Haiku 4.5} (Anthropic) as the judge, an entirely independent
model family from all evaluated models.

We sampled 100 molecules uniformly at random and generated explanations for
all 15 unique pairwise combinations of the six models, yielding
$100 \times 15 = 1{,}500$ comparisons. Win rates were aggregated per model
across all comparisons (0\% ties).

\medskip
A complete pipeline overview is illustrated in Figure~\ref{fig:pipeline}.

\begin{figure}[H]
  \centering
  % \includegraphics[width=\textwidth]{figures/pipeline.pdf}
  \fbox{\parbox{0.9\textwidth}{\centering [Pipeline figure: Molecular SMILES →
    RDKit descriptors → LightGBM → SHAP (global / local) → Prompt
    engineering → GPT-4.1-mini → NL explanation → Web app]}}
  \caption{End-to-end pipeline for explainable ADME prediction.}
  \label{fig:pipeline}
\end{figure}


%% ============================================================
%%  5. EXPERIMENTAL RESULTS
%% ============================================================
\section{Experimental Results}
\label{sec:results}

\subsection{Predictive Performance}

Table~\ref{tab:pearson} reports Pearson correlation coefficients ($r$) and
RMSE on the held-out test set for all six endpoints.

\begin{table}[H]
  \centering
  \caption{LightGBM predictive performance on the test set.}
  \label{tab:pearson}
  \begin{tabular}{lcc}
    \toprule
    \textbf{Endpoint} & \textbf{Pearson $r$} & \textbf{RMSE} \\
    \midrule
    hPPB    & 0.749 & \textit{[value]} \\
    HLM     & 0.712 & \textit{[value]} \\
    rPPB    & 0.698 & \textit{[value]} \\
    RLM     & 0.681 & \textit{[value]} \\
    MDR1-ER & 0.627 & \textit{[value]} \\
    Sol     & 0.554 & \textit{[value]} \\
    \bottomrule
  \end{tabular}
\end{table}

The highest predictive accuracy is observed for hPPB ($r = 0.749$),
consistent with the higher signal-to-noise ratio characteristic of protein
binding assays. Aqueous solubility (Sol) is the most challenging endpoint
($r = 0.554$), likely due to high experimental variability and the
multi-conformational nature of solvation.

\subsection{Global SHAP Analysis}

Figure~\ref{fig:global_shap} shows the top global SHAP features averaged
across all six endpoints. \textbf{MolLogP} (molecular lipophilicity) is
the dominant driver across all endpoints, confirming the known centrality
of lipophilicity in ADME pharmacokinetics. This finding is consistent with
König and Vellido~\cite{konig2024adme}, who reported the same leading
features using a different RDKit version, providing cross-validation of our
SHAP rankings.

\begin{figure}[H]
  \centering
  % \includegraphics[width=0.85\textwidth]{figures/global_shap.pdf}
  \fbox{\parbox{0.85\textwidth}{\centering [Global SHAP bar chart —
    MolLogP, TPSA, MolWt, NumHDonors, NumAromaticRings, \ldots]}}
  \caption{Mean absolute SHAP values for the top-10 RDKit descriptors
    (averaged across all six ADME endpoints). MolLogP is the leading driver
    in every endpoint.}
  \label{fig:global_shap}
\end{figure}

\subsection{Prompt Engineering Comparison}

Table~\ref{tab:prompt} compares the three prompt strategies for $k \in
\{2,3,5\}$ by measuring the average number of domain-relevant chemical
terms per explanation (higher = more biomedically coherent).

\begin{table}[H]
  \centering
  \caption{Average domain-specific chemical terms per explanation by
    prompt strategy and $k$ top-SHAP features.}
  \label{tab:prompt}
  \begin{tabular}{lccc}
    \toprule
    \textbf{Strategy} & \textbf{$k=2$} & \textbf{$k=3$} & \textbf{$k=5$} \\
    \midrule
    Basic        & 1.05 & 1.08 & 1.11 \\
    Biomedical   & 8.41 & \textbf{8.70} & 8.53 \\
    Reasoning    & 7.12 & 7.35 & 7.28 \\
    \bottomrule
  \end{tabular}
\end{table}

The biomedical strategy with $k=3$ features yields the highest average
domain term count (8.70), indicating greater pharmacokinetic coherence.
The basic strategy produces consistently poor domain specificity ($\approx$1
term), confirming that domain framing in the prompt is essential.

\subsection{LLM Model Comparison}

Table~\ref{tab:judge} reports pairwise win rates across the six evaluated
LLMs, as judged by Claude Haiku 4.5. GPT-4.1-mini achieves a win rate of
\textbf{96.6\%} across all 1{,}500 comparisons (100 per pair, 15 pairs),
substantially outperforming both larger open-source models and the smaller
GPT-4o-mini variant. Open-source models cluster in the 19--53\% range,
with Mistral-7B as the strongest among them.

\begin{table}[H]
  \centering
  \caption{LLM-as-a-Judge pairwise win rates (1{,}500 comparisons total,
    judge: Claude Haiku 4.5). Win rate = wins / comparisons involving that model.}
  \label{tab:judge}
  \begin{tabular}{lcc}
    \hline
    \textbf{Model} & \textbf{Wins} & \textbf{Win Rate} \\
    \hline
    GPT-4.1-mini   & 483 & \textbf{96.6\%} \\
    GPT-4o-mini    & 365 & 73.0\% \\
    Mistral-7B     & 266 & 53.2\% \\
    Gemma-7B       & 150 & 30.0\% \\
    LLaMA-3-8B     & 140 & 28.0\% \\
    Phi-3          &  96 & 19.2\% \\
    \hline
  \end{tabular}
\end{table}

\begin{figure}[H]
  \centering
  % \includegraphics[width=0.8\textwidth]{figures/model_comparison.pdf}
  \fbox{\parbox{0.8\textwidth}{\centering [Bar chart: GPT-4.1-mini 96.6\%,
    GPT-4o-mini 73.0\%, Mistral-7B 53.2\%, Gemma-7B 30.0\%,
    LLaMA-3-8B 28.0\%, Phi-3 19.2\%]}}
  \caption{Pairwise win rates (LLM-as-a-Judge, Claude Haiku 4.5, 1{,}500 comparisons).}
  \label{fig:model_comparison}
\end{figure}

\subsection{Local Explanation Example}

Figure~\ref{fig:local_example} illustrates a local explanation for a
representative molecule. The SHAP waterfall plot identifies MolLogP and
TPSA as the primary positive and negative drivers, respectively. The
GPT-4.1-mini biomedical explanation translates these values into a
coherent narrative: a high lipophilicity value is predicted to increase
hepatic clearance by promoting CYP450-mediated oxidative metabolism,
while the polar surface area slightly moderates membrane permeability.

\begin{figure}[H]
  \centering
  % \includegraphics[width=\textwidth]{figures/local_example.pdf}
  \fbox{\parbox{\textwidth}{\centering [Local explanation figure: molecule
    2D structure | SHAP waterfall | NL explanation text]}}
  \caption{Local explanation example for a sample molecule (HLM endpoint).
    Left: 2D molecular structure. Centre: SHAP waterfall plot. Right:
    GPT-4.1-mini biomedical explanation.}
  \label{fig:local_example}
\end{figure}


%% ============================================================
%%  6. DISCUSSION
%% ============================================================
\section{Discussion}
\label{sec:discussion}

\subsection{Interpretation of Results}

The dominance of MolLogP across all six ADME endpoints is consistent with
Lipinski's Rule of Five~\cite{lipinski2001experimental} and with
established pharmacokinetics: lipophilicity governs membrane permeability,
protein binding, and metabolic clearance simultaneously. This cross-endpoint
consistency reinforces the reliability of the SHAP rankings as mechanistic
signals rather than artefacts of model overfitting.

The marked superiority of the biomedical prompt strategy highlights that
LLMs require explicit domain framing to produce clinically useful
explanations. Without pharmacokinetic context, GPT-4.1-mini defaults to
general statistical language that provides little actionable guidance for
lead optimisation.

\subsection{Limitations}

\paragraph{Self-preference bias.}
The LLM-as-a-Judge evaluation uses GPT-4.1-mini as both the generator and
the judge for a subset of comparisons. Zheng~\etal~\cite{zheng2023judging}
demonstrated that LLMs exhibit systematic self-preference bias, tending to
rate outputs from the same model family more favourably. The 94.8\% win
rate for GPT-4.1-mini should therefore be interpreted cautiously; human
expert evaluation would provide a stronger validation signal.

\paragraph{Dataset scope.}
The 3{,}521-compound dataset covers drug-like chemical space but may not
generalise to highly novel scaffolds, biologics, or compounds outside the
applicability domain of the training set.

\paragraph{Descriptor coverage.}
The 217 RDKit descriptors encode global physicochemical properties but
lack 3D conformational information. Graph neural network approaches that
encode molecular topology directly may capture additional structure–property
relationships.

\paragraph{LLM hallucination.}
Although biomedical prompt constraints reduce off-topic content, LLMs can
still produce plausible-sounding but factually incorrect mechanistic claims.
Systematic red-teaming and domain-expert review of generated explanations
remain necessary before clinical deployment.

\subsection{Web Application}

The pipeline is deployed as a Streamlit application on Hugging Face Spaces,
providing open access to all six ADME models, interactive SHAP
visualisations, and GPT-4.1-mini explanations for arbitrary user-submitted
SMILES strings. This lowers the barrier for non-computational chemists and
regulatory scientists to interrogate model predictions.


%% ============================================================
%%  7. CONCLUSIONS
%% ============================================================
\section{Conclusions}
\label{sec:conclusions}

We have presented an end-to-end pipeline that combines LightGBM
regression, SHAP feature attribution, and GPT-4.1-mini generation to
produce explainable predictions of six ADME pharmacokinetic properties.
Our key findings are:

\begin{enumerate}
  \item LightGBM with RDKit descriptors achieves Pearson $r = 0.554$--$0.749$
        across six ADME endpoints, competitive with the literature baseline
        of König and Vellido~\cite{konig2024adme}.
  \item MolLogP is the dominant SHAP feature across all endpoints,
        corroborating known structure–property relationships and validating
        the explanations against the existing literature.
  \item The biomedical prompt strategy with $k=3$ SHAP features produces
        the most domain-coherent explanations (avg. 8.70 chemical terms),
        substantially outperforming basic prompting ($\approx$1.08 terms).
  \item GPT-4.1-mini achieves a 94.8\% pairwise win rate in LLM-as-a-Judge
        evaluation, though self-preference bias limits the strength of
        this conclusion.
  \item The full pipeline is openly deployed, enabling non-expert access
        to explainable ADME prediction.
\end{enumerate}

Future work will explore: (i) integration of graph neural network
descriptors for richer molecular representations; (ii) 3D-SHAP
interaction values to capture feature synergies; (iii) human expert
evaluation to complement automated LLM-as-a-Judge assessment; and
(iv) extension to additional ADME and toxicity endpoints.


%% ============================================================
%%  ACKNOWLEDGEMENTS
%% ============================================================
\section*{Acknowledgements}

The author thanks Albert Vellido and Caroline König (Universitat
Politècnica de Catalunya) for their supervision and guidance throughout
this work. % Add funding acknowledgements here if applicable.


%% ============================================================
%%  BIBLIOGRAPHY
%% ============================================================
\bibliographystyle{unsrtnat}
\bibliography{references}

% If you prefer to embed the bibliography directly (no .bib file),
% replace \bibliography{references} with \begin{thebibliography}{99} ... \end{thebibliography}

\end{document}
